\section{Every truncated polynomial thickening}
\label{sec:all-truncations}

The truncation exponent itself occurs as an exact torsion order.
On the periodic Tate complex, the curvature below is the
lifted-differential operator of
\citet[Section~1, p.~39; Proposition~1.3 and Corollary~1.4,
pp.~40--41]{Eisenbud1980}.  We require its realization as a stable
lattice endomorphism and the additive covolume that follows.
The invariant takes primary components of this curvature
endomorphism.  Their cyclic-group cohomology vanishes, so their
invariant lattices have an additive covolume, well-defined modulo
the truncation exponent.

\begin{theorem}\label{thm:all-truncations}
Let $p$ be a prime, $n\ge2$, $O=\Z_p$ or $\Z_{(p)}$, and
$A=O[e]/(e^n)$.  The degree-$p$ reduction kernel
\[
 \ker\bigl(K_0(\Lat_A(S_p))\longrightarrow
                 K_0(\Lat_O(S_p))\bigr)
\]
contains a split countably infinite direct sum of copies of $\Z/n$.
The displayed generators also vanish after inverting $p$.
In particular $\Phi_A$ is not surjective.

Let $C$ be the companion of $\Phi_p(z)$ in the monomial basis,
$a=e_{p-1}^{\mathsf T}$, and $b=e_1$.
Choose monic irreducibles $f_i\in\F_p[z]$, different from $z$,
of pairwise distinct degrees $r_i$ prime to $n(p-1)$, and integral
unit companion lifts $D_i$.  Explicit generators are
\[
 \begin{split}
 L_i&=\left(A^{pr_i},
 \begin{pmatrix}
 I_{r_i}&e(a\otimes I_{r_i})\\
 e^{n-1}(b\otimes D_i)&C\otimes I_{r_i}
 \end{pmatrix}\right),\\
 x_i&=[\operatorname{Ind}_{C_p}^{S_p}L_i]-r_i[A[S_p]].
 \end{split}
\]
The cyclic precursors $[L_i]-r_i[A[C_p]]$ likewise have exact
order $n$ and generate a split direct sum.
\end{theorem}

We first work over the complete DVR $O=\Z_p$, with fraction field
$K=\Q_p$.  A map will be called stably zero if it factors through
a finite free $A[C_p]$-lattice.  Sums and composites of such maps
have the same property.

\subsection{Curvature and integral primary components}

Let $E$ be an $A$-free $C_p$-lattice, with generator matrix $T$.
Choose any lift $\widetilde T$ over $O[[e]]$ and put
\[
 \widetilde U=\frac{\widetilde T^p-I}{e^n},\qquad
 U_E=\widetilde U\pmod{e^n}.
\]
This is an $A[C_p]$-endomorphism: the numerator is divisible
by $e^n$, and $\widetilde U$ commutes with $\widetilde T$.

\begin{lemma}\label{lem:truncation-stable-curvature}
The stable class of $U_E$ is independent of the lift and basis,
and is natural in $E$.  The groups of maps modulo stably zero maps
have characteristic $p$.
\end{lemma}

\begin{proof}
For an equivariant map $F:E\to E'$, choose lifts and write
$\widetilde T'\widetilde F-\widetilde F\widetilde T=e^n\widetilde R$.
Telescoping the difference of $p$th powers gives
\[
 U_{E'}F-FU_E
  =\sum_{i=0}^{p-1}T'^{\,p-1-i}RT^i
  =\sum_{j=0}^{p-1}T'^j H T^{-j},\qquad H=RT^{-1}.
\]
Every such group trace factors through the regular lattice:
\[
 E\longrightarrow A[C_p]\otimes_A E_{\mathrm{underlying}}
       \longrightarrow E',
\]
where the first map sends $v$ to
$\sum_jg^j\otimes T^{-j}v$, and the second sends
$g^j\otimes w$ to $T'^jH(w)$.
Changing a lift by $e^n\widetilde R$ gives the same trace formula;
a basis change is the case of an invertible $F$.
Finally the trace of an equivariant map $F$ is $pF$, proving the
last assertion.
\end{proof}

An endomorphism $U$ of a finite $O$-free module has an integral
primary decomposition
\[
 E=\bigoplus_f E_f,
\]
indexed by the irreducible factors $f\in\F_p[z]$ of its residue
characteristic polynomial.  Each $E_f/pE_f$ is $f$-primary for $U$.
This decomposition commutes with intertwined maps and is exact on
intertwined short exact sequences.  To obtain it, lift the residue
factors to $O[z]$.  A product of suitable powers of these factors
evaluated at $U$ is divisible by $p$, so a power annihilates
$E/p^kE$ for every $k$.  The factors are pairwise comaximal modulo
$p^k$.  Their Chinese remainder projectors are compatible as $k$
varies, and completeness gives the required idempotents on $E$.
Using a common annihilating product for a sequence makes the
projectors commute with its maps; projecting a chosen preimage
proves surjectivity on each component.

For $U_E$, these projectors are $A[C_p]$-linear.  Their summands are
$A$-free, since they are finite projective over the local ring $A$.
If $f\ne z$, the restriction of $U_E$ to $E_f$ is an automorphism.

\begin{lemma}\label{lem:truncation-primary-acyclic}
Every component $E_f$ with $f\ne z$ has zero Tate cohomology as an
$O[C_p]$-lattice.
\end{lemma}

\begin{proof}
We first verify that the restricted curvature can be realized by
a lift of the summand.  Choose an $A$-basis adapted to $E_f$ and
lift the change of basis to $O[[e]]$.  After conjugation, the
off-diagonal blocks of $\widetilde T$ are divisible by $e^n$.
In a diagonal block of $\widetilde T^p$, every term leaving the
summand and returning contains two such blocks, hence is divisible
by $e^{2n}$.  Thus the diagonal block $\widetilde T_f$ satisfies
\[
 \frac{\widetilde T_f^p-I}{e^n}\pmod{e^n}=U_E|_{E_f}.
\]
This curvature lift is a unit over $O[[e]]$, because its reduction
modulo the maximal ideal is a unit.

Write $\widetilde U_f$ for it.  If $(T_f-I)x=0$ modulo $e^n$,
lift $x$ and write
$(\widetilde T_f-I)x=e^ny$.  Multiplication by
$N(\widetilde T_f)$ and cancellation of $e^n$ give
\[
 \widetilde U_f x=N(\widetilde T_f)y.
\]
Therefore $x$ is a norm boundary modulo $e^n$.
If instead $N(\widetilde T_f)x=e^ny$, multiplication by
$\widetilde T_f-I$ gives
$x=(\widetilde T_f-I)\widetilde U_f^{-1}y$.
These are the two exactness assertions for the Tate complex.
\end{proof}

\subsection{The invariant lattice and its covolume}

Suppose $M$ is an $A$-free $C_p$-lattice with zero Tate cohomology.
Its invariant module $M^{C_p}$ is finite free over $O$.
After tensoring with $K$, the projector $N/p$ makes it a projective
module over the local ring $A_K=K[e]/(e^n)$, hence free of some
rank $r$.
Choose an $A_K$-basis and identify it with $A_K^r$.  The determinant
of an $O$-basis of $M^{C_p}$, relative to
$1,e,\ldots,e^{n-1}$ in these coordinates, has an integral valuation.
Define
\[
 \iota(M)=v_p(\det M^{C_p})\pmod n.
\]
This is independent of the choices.  An $O$-basis change has unit
determinant.  If $B\in\operatorname{GL}_r(A_K)$ changes the generic
basis, it induces its constant matrix $B_0$ on each of the $n$
successive $e$-adic quotients.  Consequently
\[
 \det_K(B:A_K^r\to A_K^r)=\det(B_0)^n,
\]
so the valuation changes by a multiple of $n$.

The index is additive on admissible sequences of cohomologically
trivial lattices.  Their invariants form a short exact sequence,
because the first cohomology of the left term vanishes.  Compatible
generic bases and $O$-bases give a block triangular determinant.
Also $\iota$ vanishes on finite free $A[C_p]$-lattices, whose
invariants have norm elements as an $A$-basis.

\begin{lemma}\label{lem:truncation-stable-index}
Two cohomologically trivial lattices that are stably isomorphic
have the same index $\iota$.
\end{lemma}

\begin{proof}
Let $f:M\to N$ and $g:N\to M$ be stable inverse maps, and write
$gf-I=ba$ through a projective $P$.
Then $(f,a):M\to N\oplus P$ is a split inclusion with retraction
$(g,-b)$.  Its complementary projection is
\[
 \begin{pmatrix}I-fg&fb\\-ag&I+ab\end{pmatrix}.
\]
Every entry factors through projectives: this holds for $I-fg$
by hypothesis and for the other entries because their source or
target is $P$.  The complement is therefore projective, since its
identity factors through a projective.
The ring $A[C_p]$ is finite over the local ring $A$, and is local:
its residue algebra is
$\F_p[C_p]=\F_p[g-1]/((g-1)^p)$.
Thus the complement and $P$ are free $A[C_p]$-modules.
We have an actual isomorphism after adding such modules, and their
indices vanish.
\end{proof}

\subsection{An additive primary index}

For each irreducible $f\in\F_p[z]$, $f\ne z$, set
\[
 \eta_f(E)=\iota(E_f)\in\Z/n.
\]
Lemma~\ref{lem:truncation-primary-acyclic} ensures that this is
defined.  We verify both independence and additivity.

In the stable quotient, the curvature is natural by
Lemma~\ref{lem:truncation-stable-curvature}.  On an $f$-component,
some power of $f(U)$ is divisible by $p$ and hence is stably zero.
The division by $p$ is again an equivariant $A$-linear map, since
the lattices have no $p$-torsion.  For distinct $f,g$, Bezout's
identity over $\F_p[z]$ therefore makes every stable map between
the two components zero.  An isomorphism comparing two choices of
lifts or bases, together with its inverse, becomes block diagonal
in the stable quotient.  Its diagonal blocks are stable inverses.
Lemma~\ref{lem:truncation-stable-index} proves independence of the
resulting indices.  A component occurring only on one side is
stably projective and has index zero.

For an admissible sequence, choose an $A$-linear splitting and
lift its generator matrix in upper triangular blocks.  The actual
curvature is then upper triangular with the endpoint curvatures
on its diagonal.  It is an endomorphism of the short exact sequence,
so its primary projectors give a short exact sequence
\[
 0\longrightarrow E'_f\longrightarrow E_f
       \longrightarrow E''_f\longrightarrow0.
\]
All three terms have zero Tate cohomology.  Their indices add,
and choice independence allows these compatible lifts to be used
for every original sequence.  We have proved homomorphisms
\begin{equation}\label{eq:all-truncation-index}
 \eta_f:K_0(\Lat_A(C_p))\longrightarrow\Z/n.
\end{equation}
They vanish on constant $O$-lattices, which have constant lifts
and curvature zero.  For each virtual class only finitely many
coordinates are nonzero.

\subsection{Evaluation and exact torsion orders}

\begin{proof}[Proof of Theorem~\ref{thm:all-truncations}]
For a unit companion matrix $D$, regard the displayed generator
first as a $p$ by $p$ matrix over the commutative ring
$O[D][[e]]$.  The companion cofactor identity
$a\operatorname{adj}(zI-C)b=1$ gives characteristic polynomial
\[
 (z-1)\Phi_p(z)-e^nD=z^p-1-e^nD.
\]
Cayley--Hamilton gives the exact identity
\begin{equation}\label{eq:all-truncation-companion-curvature}
 \widetilde T_D^p=I+e^n(I_p\otimes D).
\end{equation}
In particular it defines a $C_p$-lattice after reduction, and its
actual curvature is $I_p\otimes D$.
If $\overline D$ has irreducible characteristic polynomial $f$,
the entire lattice is $f$-primary.

Over $A_K$, its invariant vectors are determined by the upper
coordinate $v_0$, since
\[
 (C-I)v_c=-e^{n-1}(b\otimes D)v_0.
\]
The cokernel of $C-I$ is $O/p$, generated by the class of $b$.
Since $e^{n-1}$ retains only the constant coefficient of $v_0$,
the solution is integral precisely when that constant coefficient
is divisible by $p$.  The unit matrix $D$ does not change this
condition.  The upper row equation follows from $e^n=0$.
Projection onto $v_0$ therefore identifies the invariant lattice
with
\[
 (p,e)A^r\subset A_K^r.
\]
Its quotient in $A^r$ has $O$-length $r$.  Thus
\[
 \eta_f(L_D)=r\pmod n,
 \qquad \eta_g(L_D)=0\quad(g\ne f).
\]

Changing the generator to its $j$th power scales the curvature by
$j$.  Indeed, choosing $\widetilde T^j$ as the lift gives
\[
 \frac{\widetilde T^{jp}-I}{e^n}
 =\widetilde U\sum_{a=0}^{j-1}(\widetilde T^p)^a
 \equiv jU\pmod{e^n}.
\]
The invariant module itself is unchanged, since the two elements
generate the same cyclic group.  If $j\overline D_i$ has the same
characteristic polynomial as $\overline D_i$, comparison of
constant terms gives $j^{r_i}=1$.  The condition
$(r_i,p-1)=1$ forces $j=1$, and distinct degrees give disjoint
scalar orbits.

Mackey restriction from $S_p$ to $C_p$ gives one copy of each
normalizer twist.  Every other double coset has trivial intersection
and gives a regular cyclic lattice.  Hence
\[
 (\eta_{f_j}\circ\operatorname{Res}_{C_p}^{S_p})(x_i)
  =\begin{cases}r_i\pmod n,&i=j,\\0,&i\ne j.\end{cases}
\]
The exact sequence of Proposition~\ref{prop:truncation-torsion}
applies with the unit block $D_i$, giving $nx_i=0$ and the same
annihilator for the cyclic precursor.  The unit coordinate values
prove exact order $n$.  Multiplication of each coordinate by
$r_i^{-1}\in\Z/n$, together with finite primary support, gives
the asserted retractions onto direct sums of copies of $\Z/n$.

Reduction to $O$ gives trivial plus cyclotomic classes, with regular
class by the augmentation sequence.  After inverting $p$, the reduction is isomorphic to the regular
representation.  Lift such an isomorphism to the free modules over
$K[e]/(e^n)$, average the map over $C_p$ using $1/p$, and apply
Nakayama's lemma.  The averaged lift is an equivariant isomorphism.
This proves both kernel assertions.
Every Young lattice is constant over $O$, whereas the coordinate
invariants on the $x_i$ are nonzero.

Finally, for $O=\Z_{(p)}$, extend coefficients to
$\Z_p[e]/(e^n)$ and compose with the invariants just constructed.
This functor is exact on admissible sequences.  The explicit matrices
and their exact annihilation relations already have coefficients in
$\Z_{(p)}$, and their complete-base coordinate values are unchanged.
Thus the exact orders and the split injections descend to the stated
localized-integer base as well.
\end{proof}

\begin{corollary}\label{cor:ramified-all-truncations}
Let $O$ be an arbitrary mixed-characteristic DVR with residue
characteristic $p$, normalized valuation $v$, and absolute
ramification index $h=v(p)$.  Let $n\ge2$ and put
$A=O[e]/(e^n)$.
If $n\nmid h$, then $\Phi_A$ is not surjective in degree $p$.
If $(n,h)=1$, the degree-$p$ reduction kernel contains a split
countably infinite direct sum of copies of $\Z/n$, with generators
that also vanish after inverting $p$.
The same split assertion holds for the cyclic precursors.
No completeness, finiteness, or perfection assumption on the
residue field is required.
\end{corollary}

\begin{proof}
First assume that $O$ is complete, with uniformizer $\pi$ and
residue field $k$.  We adapt the proof of
Theorem~\ref{thm:all-truncations}, retaining only primary
polynomials $f\in\F_p[z]$, $f\ne z$.
These are \emph{coarse} primary components: if $f$ factors over $k$,
we take the sum of all the resulting primary components.  Equivalently,
reduce $f(U)$ modulo $\pi$ and take its nilpotent Fitting summand;
its complement is the invertible Fitting summand.  The characteristic
polynomial factorization over $k$, followed by the same complete-DVR
projector construction, lifts this decomposition integrally.
It is exact for intertwined sequences.  Distinct irreducibles over
$\F_p$ remain coprime over $k$, so these coarse components are disjoint,
and only finitely many occur on a fixed lattice.

The curvature trace formula is unchanged.  Stable maps are killed
by $p$, which is a unit times $\pi^h$.  On a selected component,
some power of $f(U)$ is divisible by $\pi$; its $h$th power is
therefore divisible by $p$ and is stably zero.  Division by $p$
preserves equivariance because the lattices have no $p$-torsion.
On the complement $f(U)$ is invertible.  This proves the same
stable independence of the coarse primary summand.  Since $f(0)\ne0$,
$U$ is a unit on it, and the lifted matrix-factorization argument
of Lemma~\ref{lem:truncation-primary-acyclic} proves vanishing of
its integral Tate cohomology.

Measure the invariant-lattice covolume using $v$, modulo $n$.
A change of generic $A_K$-basis still has $K$-determinant
$\det(B_0)^n$, so the ambiguity is a multiple of $n$.
Exact additivity and projective stabilization are unchanged.
We thus obtain homomorphisms
\[
 \eta_f:K_0(\Lat_A(C_p))\longrightarrow\Z/n
\]
which vanish on constant lattices and have finite primary support.
For the companion lattice $L_D$ of
Theorem~\ref{thm:all-truncations}, the curvature is again
$I_p\otimes D$, and its invariant lattice is $(p,e)A^r$.
Consequently its coarse primary coordinate is
\[
 \eta_f(L_D)=hr\pmod n.
\]

For $n\nmid h$, take $r=1$ and $D=(1)$.
Only the identity normalizer twist has coarse primary polynomial
$z-1$; the other scalar twists have distinct nonzero eigenvalues
in $\F_p$.  Mackey restriction therefore gives coordinate $h\ne0$
modulo $n$ on the induced difference from the regular lattice.
Every Young lattice is constant, so this proves non-surjectivity.

For $(n,h)=1$, choose the irreducibles $f_i$ and integral companions
$D_i$ exactly as in Theorem~\ref{thm:all-truncations}: their degrees
$r_i$ are pairwise distinct and prime to $n(p-1)$.
They exist over the fixed finite field $\F_p$, irrespective of $k$.
The scalar-orbit argument takes place over $\F_p$; its coprimality
identities remain valid after extension to $k$.
Thus the induced coordinate matrix is diagonal with entries
$hr_i\in(\Z/n)^\times$.
The same exact sequence still annihilates each displayed class
by $n$.  Normalizing the coordinates by $(hr_i)^{-1}$ gives exact
order $n$ and finite-support retractions, both cyclically and after
induction.  The reduction and generic-vanishing arguments are
unchanged.

Finally, for a noncomplete DVR, extend scalars to its completion.
This preserves the residue field and normalized valuation of $p$,
and is exact on the admissible sequences of free lattices.
Compose the complete-base invariants with this functor.  All the
explicit companion matrices and their annihilation sequences are
already defined over $O$, and their coordinate values remain $hr_i$.
The preceding conclusions therefore hold over $O$ itself.
When $(n,h)\ne1$, the value $hr$ alone does not determine the exact
order of the displayed class; no stronger torsion-order assertion
is made here.
\end{proof}

\begin{corollary}\label{cor:truncation-adams}
Let $O=\Z_{(2)}$ or $\Z_2$, $A=O[e]/(e^n)$, and
\[
 x=\left[\left(A^2,
       \begin{pmatrix}1&e\\e^{n-1}&-1\end{pmatrix}\right)\right]
       -[A[C_2]].
\]
Then $x$ has exact order $n$, $x^2=0$, and
\[
 \psi^k(x)=\begin{cases}kx,&k\text{ odd},\\0,&k\text{ even}.
                         \end{cases}
\]
For $n=3$, this gives $\psi^3(x)=0\ne x$, although $3$ is a unit
in $A$ and the cube map on $C_2$ is the identity.
\end{corollary}

\begin{proof}
Let $\sigma$ be the sign line and $P=1+\sigma$ the regular class.
For any involution lattice $(M,T)$, the maps
\[
 0\longrightarrow M\longrightarrow
 A[C_2]\otimes_A M_{\mathrm{underlying}}
 \longrightarrow\sigma M\longrightarrow0,
 \qquad
 v\longmapsto(v,Tv),\quad (u,v)\longmapsto v-Tu,
\]
form an admissible exact sequence, where the middle generator
interchanges its two coordinates.  Applied to the rank-two lattice
in the statement, it gives $[M]+[\sigma M]=2P$, and therefore
$\sigma x=-x$.
As the displayed lattice has determinant $\sigma$,
\[
 \lambda_t(x)=
 \frac{1+(P+x)t+\sigma t^2}{(1+t)(1+\sigma t)}
 =1+\frac{xt}{1-t^2}.
\]
Comparing $\lambda_t(\sigma x)=\lambda_{\sigma t}(x)$ with
$\lambda_t(-x)=\lambda_t(x)^{-1}$ gives
\[
 \left(1-\frac{xt}{1-t^2}\right)
 \left(1+\frac{xt}{1-t^2}\right)=1;
\]
its coefficient of $t^2$ is $-x^2$, so $x^2=0$.
Consequently the integral logarithmic-derivative identity becomes
\[
 \sum_{k\ge1}(-1)^{k-1}\psi^k(x)t^k
 =t\lambda_t(x)^{-1}\frac{d\lambda_t(x)}{dt}
 =\sum_{j\ge0}(2j+1)x\,t^{2j+1}.
\]
This proves the displayed values without division in the
Grothendieck ring.  Exact order follows from the theorem with
$f=z+1$.
\end{proof}

\begin{corollary}\label{cor:all-prime-two-torsion}
For every prime $p$, the degree-$p$ reduction kernel for
$O[\epsilon]/(\epsilon^2)$, with $O=\Z_{(p)}$ or $\Z_p$,
contains a split countably infinite direct sum of copies of $\Z/2$.
The displayed classes also vanish after inverting $p$.
\end{corollary}
\begin{proof}
Apply Theorem~\ref{thm:all-truncations} with $n=2$.
\end{proof}

\begin{corollary}\label{cor:adams-not-power}
Let $p$ be odd and $A=\Z_{(p)}[\epsilon]/(\epsilon^2)$.
On $K_0(\Lat_A(C_p))$ the second Adams operation defined by exterior
powers is not pullback along the automorphism $g\mapsto g^2$.
This failure occurs although two is invertible in $A$.
\end{corollary}

\begin{proof}
Exterior powers of admissible exact sequences have the usual finite
filtration.  They therefore define an additive-to-multiplicative map
$\lambda_t$ on the exact Grothendieck group, and
$\lambda_t(-x)=\lambda_t(x)^{-1}$.
Let $x=[L_1]-[A[C_p]]$ be the nonzero order-two class of
Proposition~\ref{prop:truncation-torsion}, with $n=2$.
As $-x=x$, comparison of coefficients of $t^2$ gives
\[
 \lambda^2(x)=x^2-\lambda^2(x),\qquad
 \psi^2(x)=x^2-2\lambda^2(x)=0.
\]
Pullback along $g\mapsto g^2$ is an automorphism of the exact category,
so its induced group automorphism cannot kill $x$.
\end{proof}

Thus a formula on ordinary characters cannot determine the Adams
operations on this integral Grothendieck group.  The covolume argument supplies the independent information needed
to determine the truncation classes.
The annihilation identity used in the proof is the degree-two instance
of the general torsion--Adams relation of
\citet[Proposition~2, p.~2]{Clauwens2010}; the proof above derives the
needed identity directly from the exterior-power addition formula.

In the projective group-ring setting,
\citet[Propositions~4.2.3 and~4.3.2, pp.~1490--1492]{Pappas2015}
compares Adams operations with Cassou-Nogu\`es--Taylor operations.
The examples here concern lattices free over singular coefficient
rings, whose exact classes need not be determined by ordinary characters.
